\SourceLeaf{10}
\[
\mathrm{d}\Sigma=-\mathrm{d}s'.\cos QX'=+\mathrm{d}s''.\cos QX''=-\mathrm{d}s'''.\cos QX'''\text{ etc.}
\]
adeoque propter partium multitudinem parem
\[
\mathrm{d}s'.\cos QX'+\mathrm{d}s''.\cos QX''+\mathrm{d}s'''.\cos QX'''+\text{ etc.}=0
\]
\par
Tractando eodem modo omnia reliqua elementa $\mathrm{d}\Sigma$, atque summando, nanciscimur
\par\medskip\Needspace{4\baselineskip}
\begin{center}THEOREMA PRIMUM.\end{center}
\par
\textit{Integrale $\int\mathrm{d}s\cos QX$ per totam corporis superficiem extensum fit $=0$.}
\par
Generalius eodem modo invenitur, integrale
\[
\int(T\cos QX+U\cos QY+V\cos QZ)\mathrm{d}s
\]
evanescere, si $T,U,V$ resp. designent functiones rationales solarum $y,z$ solarum $x,z$ solarumque $x,y$.

\Article{4}
\par
Quum volumina partium cylindri a plano nostro usque ad puncta $P',P'',P'''$ etc resp. sint $=\mathrm{d}\Sigma.(x'-\alpha)$, $\mathrm{d}\Sigma.(x''-\alpha)$, $\mathrm{d}\Sigma.(x'''-\alpha)$ etc., pars voluminis corporis ea, quae intra cylindrum sita est, erit
\[
\begin{aligned}
&=-x'\mathrm{d}\Sigma+x''\mathrm{d}\Sigma-x'''\mathrm{d}\Sigma+\text{ etc.}\\
&=\mathrm{d}s'.x'\cos QX'+\mathrm{d}s''.x''\cos QX''+\mathrm{d}s'''.x'''\cos QX'''+\text{ etc.}
\end{aligned}
\]
unde summando pro omnibus $\mathrm{d}\Sigma$ obtinemus
\par\medskip\Needspace{4\baselineskip}
\begin{center}THEOREMA SECUNDUM.\end{center}
\par
\textit{Volumen integrum corporis exprimitur per integrale $\int\mathrm{d}s.x\cos QX$ per totam superficiem extensum.}
\par
Manifesto idem volumen etiam per $\int\mathrm{d}s.y\cos QY$ vel per $\int\mathrm{d}s.z\cos QZ$ exprimere licebit.

\Article{5}
\par
Concipiatur iam primo cylinder totus materia uniformiter densa repletus, videamusque quantam singula eius elementa attractionem in punctum $M$ exer\SourceLeaf{11}ceant. Dividatur cylinder per plana infinite sibi proxima basique parallela in cylindros elementares, qualium unus, ad punctum cuius coordinatae sunt $\xi,\eta,\zeta$, per $\mathrm{d}\Sigma.\mathrm{d}\xi$ exprimi poterit. Huius distantia a puncto $M$ erit
\[
=\BareRoot{((a-\xi)^2+(b-\eta)^2+(c-\zeta)^2)}=\rho
\]
unde ipsius attractio in punctum $M$ exhiberi poterit per $\mathrm{d}\Sigma.\mathrm{d}\xi.f\rho$, denotante functione $f\rho$ legem attractionis. Quare quum per totum cylindrum sola $\xi$ tamquam variabilis spectanda sit, erit $\rho\mathrm{d}\rho=-(a-\xi)\mathrm{d}\xi$, et proin attractio elementi $=-\frac{\rho f\rho.\mathrm{d}\rho.\mathrm{d}\Sigma}{(a-\xi)}$. Qua resoluta in tres attractiones partiales axibus coordinatarum $x,y,z$ parallelas atque oppositas, prima erit $=-f\rho.\mathrm{d}\rho.\mathrm{d}\Sigma$. Hinc designando integrale $\int f\rho.\mathrm{d}\rho$ per $F\rho$, attractio cylindri a basi $\mathrm{d}\Sigma$ usque ad punctum cuius coordinata prima $=\xi$ in punctum $M$ secundum axem coordinatarum $x$ erit $=-(F\rho-\mathrm{Const.})\mathrm{d}\Sigma=-(F\rho-FR)\mathrm{d}\Sigma$, si $R$ supponitur designare distantiam basis $\mathrm{d}\Sigma$ a puncto $M$. Hinc sequitur, eandem attractionem partialem omnium partium corporis, quae intra cylindrum iacent, fieri
\[
\begin{aligned}
&=(Fr'-Fr''+Fr'''-\text{ etc.})\mathrm{d}\Sigma\\
&=-Fr'.\mathrm{d}s'.\cos QX'-Fr''.\mathrm{d}s''.\cos QX''-Fr'''.\mathrm{d}s'''.\cos QX'''-\text{ etc.}
\end{aligned}
\]
\par
Extendendo haec ratiocinia ad omnia elementa $\mathrm{d}\Sigma$, colligimus
\par\medskip\Needspace{4\baselineskip}
\begin{center}THEOREMA TERTIUM.\end{center}
\par
\textit{Attractio corporis in punctum $M$, axi coordinatarum $x$ parallela atque opposita, exhibetur per integrale $-\int Fr.\mathrm{d}s.\cos QX$ per totam superficiem extensum.}
\par
{\rightskip=0pt plus 3em Prorsus simili modo manifesto attractio secundum duas reliquas directiones principales exprimetur per integralia $-\int Fr.\mathrm{d}s.\cos QY$, $-\int Fr.\mathrm{d}s.\cos QZ$.\par}

\Article{6}
\par
Iam rem alia via aggrediemur. Concipiatur superficies sphaerica radio $=1$ circa centrum $M$ descripta, atque in elementa infinite parva dispertita. Sit $\Pi$ punctum huius superficiei ad spatiolum $\mathrm{d}\Sigma$ in eadem pertinens; ducatur radius $M\Pi$, atque si opus est ultra sphaerae superficiem indefinite producatur. Sint $P',P'',P''$ etc. puncta, in quibus hic radius superficiem corporis nostri deinceps secat, excluso tamen ipso puncto $M$, si forte in ipsa superficie iacet. Horum
